\documentclass[11pt]{article}

\usepackage[a4paper,margin=1in]{geometry}
\usepackage{amsmath,amssymb}
\usepackage{bm}
\usepackage{booktabs}
\usepackage{hyperref}
\usepackage{enumitem}

\title{A Physics-Informed Neural Surrogate for Stellar Chemical Equilibrium in \texttt{jorg}}
\author{Implementation Design Note}
\date{\today}

\begin{document}
\maketitle

\begin{abstract}
This document specifies the current PINN design for fast chemical-equilibrium inference in \texttt{jorg}. The reference physics follows the Korg formulation implemented in \path{Korg.jl/src/statmech.jl} and ported in \path{jorg/src/jorg/statmech/korg_chemical_equilibrium.py}. The PINN predicts electron density and per-element neutral fractions from thermodynamic state and abundances, and is trained with a two-stage objective: (i) supervised regression to trusted solver outputs, and (ii) residual minimization of atomic element conservation and charge neutrality. The present physics loss is atomic-only; molecular constraints are delegated to the label generator and not yet enforced explicitly in the PINN objective.
\end{abstract}

\section{Problem Statement and Scope}
Given
\begin{align}
    \left(T,\ n_{\mathrm{tot}},\ \bm{A}\right), \qquad
    \bm{A} = (A_1,\dots,A_{92}), \quad \sum_{Z=1}^{92} A_Z = 1,
\end{align}
we seek a fast approximation to the equilibrium solution
\begin{align}
    \left(n_e,\ \bm{f}\right), \qquad \bm{f} = (f_1,\dots,f_{92}),
\end{align}
where $n_e$ is electron number density and $f_Z$ is the neutral fraction of element $Z$.

The reference equilibrium labels are produced by the trusted backends:
\begin{itemize}[leftmargin=1.5em]
    \item \texttt{jorg/src/jorg/statmech/chem\_eq\_jax.py} (default),
    \item \texttt{jorg/src/jorg/statmech/korg\_chemical\_equilibrium.py} (SciPy port of Korg).
\end{itemize}
The current PINN loss enforces only atomic constraints (neutral, singly ionized, doubly ionized species), consistent with the present phase of development.

\section{Reference Physics: Atomic Ionization Structure}
\subsection{Saha ionization weights}
For each element $Z$, define Saha ratios
\begin{align}
    w_{Z,\mathrm{II}} &= \frac{n_{Z,\mathrm{II}}}{n_{Z,\mathrm{I}}}, \\
    w_{Z,\mathrm{III}} &= \frac{n_{Z,\mathrm{III}}}{n_{Z,\mathrm{I}}}.
\end{align}
Following \texttt{statmech.jl}, the first-ionization ratio is
\begin{align}
    w_{Z,\mathrm{II}} =
    \frac{2}{n_e}
    \frac{U_{Z,\mathrm{II}}(T)}{U_{Z,\mathrm{I}}(T)}
    \left(\frac{2\pi m_e k_B T}{h^2}\right)^{3/2}
    \exp\!\left(-\frac{\chi_{Z,1}}{k_B T}\right),
\end{align}
and the second-ionization ratio is
\begin{align}
    w_{Z,\mathrm{III}} =
    w_{Z,\mathrm{II}}
    \frac{2}{n_e}
    \frac{U_{Z,\mathrm{III}}(T)}{U_{Z,\mathrm{II}}(T)}
    \left(\frac{2\pi m_e k_B T}{h^2}\right)^{3/2}
    \exp\!\left(-\frac{\chi_{Z,2}}{k_B T}\right),
\end{align}
with the hydrogen special case $w_{1,\mathrm{III}}=0$ in the implementation.

\subsection{Atomic number densities}
Define the nucleus reservoir for element $Z$ as
\begin{align}
    n_{\mathrm{nuc},Z} = A_Z\,(n_{\mathrm{tot}} - n_e).
\end{align}
Given neutral fraction $f_Z$, the three atomic stages used in the PINN loss are
\begin{align}
    n_{Z,\mathrm{I}} &= f_Z\, n_{\mathrm{nuc},Z}, \\
    n_{Z,\mathrm{II}} &= w_{Z,\mathrm{II}}\, n_{Z,\mathrm{I}}, \\
    n_{Z,\mathrm{III}} &= w_{Z,\mathrm{III}}\, n_{Z,\mathrm{I}}.
\end{align}

\section{PINN Parameterization}
\subsection{Inputs and outputs}
The model input is the 94-dimensional vector
\begin{align}
    \bm{x} = \left[\log_{10}T,\ \log_{10}n_{\mathrm{tot}},\ \bm{A}\right].
\end{align}
The network output has 93 channels:
\begin{align}
    (\bm{z}_f,\ z_e), \qquad \bm{z}_f\in\mathbb{R}^{92},\ z_e\in\mathbb{R}.
\end{align}
It is mapped to physical variables as
\begin{align}
    \hat{f}_Z &= \mathrm{clip}\!\left(\sigma(z_{f,Z}),\ 10^{-12},\ 1\right), \\
    \hat{n}_e &= n_{\mathrm{tot}}\,\sigma(z_e),
\end{align}
which guarantees $0<\hat{f}_Z\le 1$ and $0<\hat{n}_e<n_{\mathrm{tot}}$.

\subsection{Architecture}
The baseline architecture in \texttt{chem\_eq\_pinn\_models.py} is an MLP with hidden widths
\begin{align}
    (256,\ 512,\ 512,\ 256,\ 128)
\end{align}
and Swish activation. A Flax implementation is used when available, with a pure-JAX fallback path preserving the same functional mapping.

\section{Training Objectives}
\subsection{Supervised objective}
Training labels $(n_e^\star,\bm{f}^\star)$ are generated from the reference solvers. The supervised objective is
\begin{align}
    \mathcal{L}_{\mathrm{sup}}
    &= \mathcal{L}_{n_e} + \mathcal{L}_f, \\
    \mathcal{L}_{n_e}
    &= \left\langle \left[\log_{10}\hat{n}_e - \log_{10}n_e^\star\right]^2 \right\rangle, \\
    \mathcal{L}_f
    &= \left\langle \left[\log_{10}\hat{f}_Z - \log_{10}f_Z^\star\right]^2 \right\rangle,
\end{align}
with numerical floors applied in code before logarithms.

\subsection{Physics objective}
The PINN physics residuals are computed from $\hat{n}_e$ and $\hat{\bm{f}}$.

\paragraph{Element conservation residual}
\begin{align}
    r_{\mathrm{elem},Z}
    = \frac{
        n_{Z,\mathrm{I}} + n_{Z,\mathrm{II}} + n_{Z,\mathrm{III}} - n_{\mathrm{nuc},Z}
    }{
        \max(n_{\mathrm{nuc},Z},\epsilon)
    },
    \qquad \epsilon = 10^{-100}.
\end{align}
The loss uses abundance-weighted MSE:
\begin{align}
    \mathcal{L}_{\mathrm{elem}}
    = \left\langle
    \frac{1}{92}\sum_{Z=1}^{92}\omega_Z\,r_{\mathrm{elem},Z}^2
    \right\rangle, \qquad
    \omega_Z = \frac{\max(A_Z,10^{-12})}{\langle \max(A_Z,10^{-12})\rangle_Z}.
\end{align}

\paragraph{Charge-neutrality residual}
\begin{align}
    r_{\mathrm{chg}}
    = \frac{
        \sum_{Z=1}^{92}\left(n_{Z,\mathrm{II}}+2n_{Z,\mathrm{III}}\right)-n_e
    }{
        \max(n_e,\epsilon)
    },
    \qquad
    \mathcal{L}_{\mathrm{chg}}=\langle r_{\mathrm{chg}}^2\rangle.
\end{align}

The total physics loss is
\begin{align}
    \mathcal{L}_{\mathrm{phys}}
    = w_{\mathrm{elem}}\mathcal{L}_{\mathrm{elem}} + w_{\mathrm{chg}}\mathcal{L}_{\mathrm{chg}},
\end{align}
with current defaults $w_{\mathrm{elem}}=1$, $w_{\mathrm{chg}}=10$.

\subsection{Two-phase optimization}
The default schedule is:
\begin{enumerate}[leftmargin=1.5em]
    \item supervised pretraining with $\mathcal{L}_{\mathrm{sup}}$;
    \item phase-2 fine tuning with either $\mathcal{L}_{\mathrm{phys}}$ only or
    \begin{align}
        \mathcal{L}_{\mathrm{phase2}}
        = w_{\mathrm{sup,2}}\mathcal{L}_{\mathrm{sup}} + \mathcal{L}_{\mathrm{phys}}.
    \end{align}
\end{enumerate}
Optimization uses AdamW with warmup and cosine decay (\texttt{chem\_eq\_pinn\_train.py}).

\section{Data Generation and Solver Coupling}
Training/validation tuples are
\begin{align}
    \left(T,\ n_{\mathrm{tot}},\ \bm{A},\ n_e^\star,\ \bm{f}^\star\right),
\end{align}
sampled over broad stellar-atmosphere ranges in $(T,\log_{10}n_{\mathrm{tot}},[M/H],[\alpha/H])$ and filtered for physically valid solutions.

Two points are important for interpretation:
\begin{itemize}[leftmargin=1.5em]
    \item The reference solvers may include molecular equilibrium terms, depending on supplied equilibrium-constant tables.
    \item The current PINN physics loss does not yet include explicit molecular residual equations.
\end{itemize}
Therefore, molecular effects are presently learned through supervision rather than enforced as hard residual constraints.

\section{Software Mapping}
\begin{table}[h!]
\centering
\begin{tabular}{p{0.43\linewidth}p{0.5\linewidth}}
\toprule
File & Functional role \\
\midrule
\path{chem_eq_pinn_models.py} & Network architecture and output parameterization \\
\path{chem_eq_pinn_loss.py} & Saha weights, residuals, and physics losses \\
\path{chem_eq_pinn_train.py} & Data generation, objectives, and optimization loops \\
\path{chem_eq_pinn.py} & High-level train/infer/save/load API \\
\path{korg_chemical_equilibrium.py} & Full SciPy reference solver (Korg port) \\
\path{Korg.jl/src/statmech.jl} & Upstream physical formulation and algorithmic reference \\
\bottomrule
\end{tabular}
\end{table}

\section{Validation Protocol}
Model quality is assessed with:
\begin{itemize}[leftmargin=1.5em]
    \item supervised agreement metrics (e.g. mean and percentile relative errors in $n_e$ and $f_Z$),
    \item physics diagnostics ($|r_{\mathrm{elem},Z}|$ and $|r_{\mathrm{chg}}|$ distributions),
    \item out-of-sample checks over temperature, density, and metallicity slices.
\end{itemize}
For operational reproducibility, cached datasets are recommended via
\path{examples/train_chem_eq_pinn_full.py} options:
\begin{itemize}[leftmargin=1.5em]
    \item \texttt{--prepare-data-only},
    \item \texttt{--load-data},
    \item \texttt{--data-dir}.
\end{itemize}

\section{Current Limitations and Next Steps}
\begin{itemize}[leftmargin=1.5em]
    \item Physics regularization is atomic-only; molecular equilibrium constraints are not yet explicit in $\mathcal{L}_{\mathrm{phys}}$.
    \item The surrogate is calibrated over a finite thermodynamic domain; extrapolation to extreme conditions requires dedicated stress testing.
    \item A direct extension path is to augment the residual system with molecular conservation equations consistent with the Korg formalism.
\end{itemize}

\end{document}
